% Generated draft: requires mathematical and source review.
\documentclass[11pt]{article}
\usepackage[margin=25mm]{geometry}
\usepackage{fontspec}
\setmainfont{Latin Modern Roman}
\usepackage{amsmath,amssymb,mathtools,unicode-math}
\setmathfont{Latin Modern Math}
\usepackage{enumitem,needspace,xurl,hyperref}
\setlength{\parindent}{0pt}
\setlength{\parskip}{5pt}
\setlength{\emergencystretch}{3em}
\newcommand{\N}{\mathbb{N}}
\newcommand{\Z}{\mathbb{Z}}
\newcommand{\Q}{\mathbb{Q}}
\newcommand{\R}{\mathbb{R}}
\newcommand{\C}{\mathbb{C}}
\newcommand{\F}{\mathbb{F}}
\newcommand{\A}{\mathbb{A}}
\newcommand{\PP}{\mathbb{P}}
\newcommand{\E}{\mathbb{E}}
\newcommand{\eps}{\varepsilon}
\newcommand{\dd}{\,\mathrm{d}}
\begin{document}
\section*{UnsolvedMath ID: 20000807}
\textbf{Source problem code:} catalogue entry
\subsection{Definitions and mathematical statement}
Let \(G_\Q=\operatorname{Gal}(\overline\Q/\Q)\). The Hilbert scheme \(\operatorname{Hilb}^d(\A^N_\Q)\) parametrizes length-\(d\) finite subschemes. Its base change to \(\overline\Q\) has geometric irreducible components on which \(G_\Q\) acts. The first original question asks for explicit \(d,N\) and a geometric component \(C\) such that \(\sigma(C)\neq C\) for some \(\sigma\in G_\Q\), or a proof this cannot occur. The second asks the same for the Hilbert schemes \(\operatorname{Hilb}^{P}(\PP^3)\) with a curve Hilbert polynomial \(P\). The extracted \(\operatorname{Gal}(\Q/\Q)\) is read with the missing algebraic-closure bar restored.

The supplied existence theorem for the first clause is stronger in its Galois quantifier:
\[
 \forall\sigma\in G_\Q\setminus\{1\}\ \exists d,N\geq1\ \exists C\subset\operatorname{Hilb}^d(\A^N_{\overline\Q})
 \text{ a geometric irreducible component}:\quad\sigma(C)\neq C.
\]
The Hilbert scheme may depend on \(\sigma\); no single finite-type Hilbert scheme is claimed to carry a faithful action on its finite set of components. The aid derives this by combining Easton--Vakil's faithful Galois action on surface-moduli components with Jelisiejew's functorial truncation and TNT-frame section--retraction constructions. It reports transfer of nontrivial component constant fields at a generic point, followed by spreading of the frame. If the source surface is canonically embedded in \(\PP^r\), the construction gives \(N=2r+8\) and length \(d=\dim_k(T/J)\) for the resulting finite frame algebra over its coefficient field \(k\). The actual \(r,T,J\), field, and numerical \(d\) are not supplied. This is a constructive existence route, not the compact fully explicit example requested by the source.

The aid's independent degeneration statement is: every geometric irreducible component in the point-Hilbert schemes above, and in the \emph{full} projective Hilbert scheme for a fixed polynomial, contains a monomial-ideal point defined over \(\Q\). Consequently if \(C\) is moved, there is one \(\Q\)-rational monomial point \(m\in C\) such that
\[
 m\in\gamma(C)\qquad\text{for every }\gamma\in G_\Q.
\]
Thus all conjugates meet at \(m\); when two of them are distinct, \(m\) is singular on the ambient Hilbert scheme. This does not assert that \(m\) is intrinsically singular on each individual component, or that a rational point forces the component to descend to \(\Q\).

The reported proof uses monomial degeneration and the fact that a connected torus preserves each irreducible component, followed by the rationality of monomial ideals. For projective curves, the statement concerns the full Hilbert scheme, allowing the scheme-theoretic degenerations with the same polynomial; it is not silently asserted inside an open locus of smooth, integral, or locally Cohen--Macaulay curves that might exclude the monomial limit. The source does not specify which curve convention it intends.

The report found no answer to the second clause in fixed \(\PP^3\). This is a limitation of the supplied work, not a proof of nonexistence or a verified exhaustive literature conclusion. The first-clause existence transfer and the monomial collision condition do not themselves produce a Galois-moved curve component in that fixed projective space.
\subsection{Short English statement}
The supplied construction yields a moved point-Hilbert component for every nontrivial Galois element, but does not provide small numerical parameters or explicit defining data. A separate result says each component contains a rational monomial degeneration, so a moved component meets all its conjugates at an ambient singular point. The fixed-\(\PP^3\) curve question is not answered in the report.
\subsection{Sources}
\begin{itemize}
\item Local registry: permanent UnsolvedMath ID \(20000807\), exact public code \texttt{catalogue entry}; \texttt{aim-arithmetic-geometry-notes.json}, \texttt{notes[52]}.
\item AIM workshop \emph{Components of Hilbert schemes}, Problem 6, clauses (1) and (2): \url{https://aimath.org/WWN/hilbertschemes/hilbertschemes.pdf}. The supplied URL is provenance; no PDF access is claimed.
\item Supplied research summary, ``Galois-moving point-Hilbert components via constant-field transfer,'' attempt 1, cites Easton--Vakil's surface-moduli component action and Jelisiejew's functorial truncation/TNT-frame retractions without full paper titles, dates, or theorem citations. Original remarks and literature fields are empty.
\end{itemize}
\end{document}
